\subsection{Trunk Static Optimization With a Selectable Cost Exponent}
\label{sec:spn-10b}

This section documents \texttt{spine/static\_optimization.py} and the shared
numeric core \texttt{static\_opt\_core.py} (issue \#12761, epic \#12735). It
resolves trunk muscle redundancy against the spinal inverse-dynamics
generalized forces of Section~\ref{sec:spn-10a}, with the activation cost
$\sum a^p$, $p\in\{2,3\}$, and declared reserve actuators per spinal degree of
freedom.

\paragraph{Inputs and Frames.}
Static optimization needs the generalized force $\tau_j$ of every level
coordinate $j$, not only the per-level moments in the inferior ISB frame.
\texttt{spinal\_generalized\_forces} reuses the same ID solve as
\texttt{run\_spinal\_inverse\_dynamics} (one shared private pass, no second
solver) and returns $q$, $\dot q$ and $\tau$ per model coordinate. With
\texttt{exclude\_muscles} (the default) every muscle is switched off in the ID
state, so $\tau$ is the net generalized force the muscles and reserves must
supply; a muscle-bearing model then gives the same $\tau$ as its muscle-free
skeleton (tested to $10^{-9}$). Couplers prescribe kinematics only and are
disabled in both the ID and the static-optimization model copies (ADR-0055
item 6). Generalized forces are in N$\cdot$m and follow each coordinate's
OpenSim sign; muscle forces are tendon tensions in N.

\paragraph{Muscle Model.}
Each muscle is evaluated natively with a rigid tendon
(\texttt{ignore\_tendon\_compliance}): fibre length and velocity follow from
the path length and lengthening speed, and the tendon force is affine in the
activation,
\[
  F_m(a_m) = P_m + a_m A_m, \qquad A_m = F_{0,m} f_L f_V \cos\alpha ,
\]
with $P_m$ the passive and fibre-damping force from OpenSim's own Millard
curves. $A_m$ and $P_m$ are read at $a=0.5$ and $a=1$ and extrapolated; this is
exact above the Millard minimum activation (0.01 by default, below which
OpenSim clamps). Activation dynamics and tendon compliance are not represented.
Moment arms $R_{mj}$ come from \texttt{GeometryPath.computeMomentArm} at the
frame pose. The muscle generalized force is $M(a) = R^{\mathsf T}(P + a\circ A)
= o + G a$.

\paragraph{Optimization Problem.}
With declared reserve limits $L_j>0$ the frame problem is lexicographic. Stage
one finds the smallest weighted remainder
\[
  r^\ast = \tau - o - G a_1, \qquad
  a_1 \in \arg\min_{0\le a\le 1} \sum_j \bigl((\tau_j - o_j - (Ga)_j)/L_j\bigr)^2 ,
\]
a bounded linear least-squares problem (BVLS) whose remainder $r^\ast$ is
unique even when $a_1$ is not. Stage two resolves the redundancy,
\[
  \min_{0\le a\le 1} \sum_m a_m^p \quad \text{s.t.}\quad G a = G a_1 ,
\]
by SLSQP with the analytic gradient $p\,a^{p-1}$ and the constant constraint
Jacobian, the equality reduced to an orthonormal basis of the column space of
$G$ so that it has full row rank. For $p=2$ stage two is a convex quadratic
program (SLSQP is exact on it); for $p=3$ it is a smooth convex program on
$a\ge0$. A final minimum-norm Newton step on the unbounded activations removes
residual constraint error. This is the $w\to\infty$ limit of the penalty
$\sum a^p + w\sum (r_j/L_j)^2$: reserves carry exactly what the muscles cannot
produce and are zero to solver precision when $\tau$ is reachable. The leg
static optimization (\texttt{musculoskeletal\_static\_opt}, issue \#11617)
keeps its OpenSim-style $L_2$ penalty \texttt{solve\_frame}, which moved
unchanged into the shared core (golden-tested to $10^{-12}$).

\paragraph{Reserve Reporting and Flags.}
Per frame the result reports activations, muscle forces, the muscle part
$M(a)$, the reserves $r$ (with $M(a)+r=\tau$ to $10^{-9}$) and the reserve
share $s = \sum_j |r_j| / \sum_j |\tau_j|$. A frame is flagged, never dropped,
when $s$ exceeds the declared \texttt{max\_share} (default 0.05), when any
$|r_j|>L_j$ (default 1 N$\cdot$m), or when an optimizer stage reports failure.
The activation participation ratio $(\sum a)^2/\sum a^2$ measures how many
muscles share the load.

\paragraph{Runtime Contract.}
Frames are selected at named events (address, top and impact, nearest sample)
plus every \texttt{stride}-th sample; a selection larger than
\texttt{max\_frames} is an error, not a silent truncation. Continuous-integration
runs use the fixture with at most five frames.

\paragraph{Fixture Design.}
The eight lumped bilateral fixture muscles of Section SPN-7b cannot span the 15
rotational degrees of freedom of the five-level fixture, so its ID moments
are in general unreachable. The tests therefore use a one-level variant
(level L3--L4 only: the muscles run pelvis to L3 to the thorax proxy and only
the three L3--L4 coordinates lie between pelvis and L3), where eight muscles
act on three degrees of freedom and the gravity ID moments are reachable with
redundancy five, and the five-level fixture, where the reserves must satisfy
the stage-one optimality (KKT) conditions and every frame is flagged
(measured share 0.38 to 0.46). All values are invented.

\paragraph{Acceptance Evidence.}
On the one-level fixture, for $p=2$ and $p=3$, the muscle part meets the ID
generalized forces (about 23 N$\cdot$m) to $10^{-6}$ with reserves below
$10^{-6}$ (measured $10^{-14}$) at the three event frames; a symmetric pose
gives left and right activations equal to $10^{-6}$; and $p=3$ spreads load
across more muscles than $p=2$ in every frame (participation ratio 4.9 to 5.7
against 4.4 to 5.0). The analytic gradient matches central finite differences.

\paragraph{Early-Sample Failure: Diagnosis and Fix.}
Issue \#12827. On the five-sample synthetic study motion (times 0 to 0.4 s,
invented fixture of Section SPN-13a, reserve limits $10^{6}$ N$\cdot$m, $p=2$
and $p=3$) the solver reported failure at the first two samples ("Inequality
constraints incompatible") and success at the other three. Per-frame record
before the fix: the demands are 20 to 30 N$\cdot$m with a weighted reserve
remainder of 3.5 to 5.2 N$\cdot$m (the eight muscles cannot reach 15
degrees of freedom, as designed); stage one (bounded least squares) succeeds in
every frame; stage two (SLSQP) stops after one iteration at the stage-one point
in frames~0 and~1, with the equality-constraint violation already $10^{-16}$,
all activations below 0.16 (none at the upper bound), five of eight at zero.
Hypotheses, tested one at a time:
\begin{enumerate}
  \item \emph{End-sample derivative error (rejected).} The motion was
    resampled at 5, 9, 21 and 81 samples and the generalized forces and speeds
    at $t=0$ and $t=0.1$ s compared with the 81-sample run. The not-a-knot
    spline error is real but small: maximum $|\Delta\tau|$ of 0.053, 0.015 and
    0.0023 N$\cdot$m and $|\Delta\dot q|$ of $3.3\times10^{-4}$,
    $4.6\times10^{-5}$ and $3.0\times10^{-6}$ rad/s for 5, 9 and 21 samples
    (at $t=0.1$ s: 0.0066, 0.0027, 0.00025 N$\cdot$m). The failure persists at
    81 samples, where the error is zero by construction, and at the same two
    times, so derivative quality is not the cause.
  \item \emph{Initial guess (rejected).} Restarting stage two from
    $0.5\,a_1$, from $a_1$ plus $10^{-3}$ and from the all-0.5 point gave the
    same status in every case but one (all-0.5 at the 81-sample $t=0.08$ s
    frame converged), so the start point is not the primary cause.
  \item \emph{Activation bounds and reserve sizing (rejected).} The reserve
    limit was set to 1 and to $10^{6}$ N$\cdot$m: identical failures, identical
    activations. No activation reaches 0.16 in the failing frames, so the
    upper bound is inactive.
  \item \emph{Solver options (reported, not changed).} Tightening or relaxing
    the SLSQP function tolerance ($10^{-14}$, $10^{-10}$, $10^{-8}$) does not
    change the failure. No tolerance was altered in the fix.
  \item \emph{A muscle with no force capacity (confirmed).} In the failing
    frames the active-force gain of \texttt{rectus\_abdominis\_r} is exactly
    zero. Natively, its path length minus the tendon slack length (0.0476 m
    at the first sample) is below the minimum fibre length (0.0533 m), so the
    rigid-tendon Millard muscle is clamped and its force is zero at every
    activation; the length grows with the motion and the muscle develops
    21.8, 129.7 and 212.9 N per unit activation at samples~3 to~5. The zero
    column of the activation gain breaks SLSQP: on 150 random invented
    problems with one all-zero column, 98 failed and every one of them
    succeeded with the column removed; giving the column a gain of 0.001 (or
    1) in the fixture frames also cured it, while $10^{-6}$ did not at
    sample~1.
\end{enumerate}
\emph{Fix.} A muscle whose gain column is identically zero cannot change the
muscle generalized force and its cost is minimal at zero, so
\texttt{static\_opt\_core.\_redistribute} holds it at zero activation and
leaves it out of the stage-two problem; if no muscle has capacity the
activations are zero and the whole demand is a reserve. The optimum is
unchanged (the muscle contributes nothing either way); only the degenerate
variable is removed. No tolerance, budget or option changed. After the fix all
five samples succeed for $p=2$ and $p=3$ with the muscle at zero activation in
the first two. Near-zero (not exactly zero) capacities are not handled
specially; the $10^{-6}$ experiment above shows that fragility remains
possible for capacities of that order, which the fixture does not produce.
The reserve share in these frames stays high by design of the fixture and the
frames remain \texttt{unqualified}.

\paragraph{Native Bruno Run on Capture-A.}
Static optimization was evaluated on 2026-10-19 on main \texttt{fc49cd6e6f}
under OpenSim 4.6 using the licensed model \texttt{bruno2015\_fullbody\_os4\_v2}
(model SHA-256
\texttt{e488ae2cf67be3de36f569cab514e181bfa89d5186d998f1627b461b4498fd47})
and capture-A continuous IK (SHA-256
\texttt{69fd5b8ea56b237bbbcd7e733a72807d087c831913fbf9de59f09c6ed4e8fc93};
530 frames, 1.469\,s). Body weight is 764.6\,N, 446 Thelen2003 muscles are
swapped for Millard2012 (\texttt{millard\_equivalent}), 22 thoracic LB/AR
coordinates are unlocked, and coupling law is \texttt{bruno2015\_regional}.
The solver is lexicographic: bounded linear least squares (BVLS) reserve
minimization in stage one, followed by $\min \sum a^p$ in stage two ($p=3$ via
\texttt{clarabel\_irls}; $p=2$ via \texttt{clarabel\_qp}). Across 51 DOFs (17
levels from L5--S1 to T1--T2 $\times$ 3 axes) with reserve limit 1\,N$\cdot$m per
DOF and max share 0.05, all 3 event frames solved (0 unsuccessful) and all 3 are
flagged for reserve use. The maximum $|r_j|$ is 341.6\,N$\cdot$m. All stages
completed with status \texttt{ok}, but the receipt is \texttt{unqualified}
(\texttt{scientific\_status unqualified}) due to identity torso alignment, pelvis
residuals exceeding budget, and the 3 flagged frames.

\begin{center}
\begin{tabular}{lrrr}
\hline
Event & Reserve Share & Activations ($a=1$ / $a=0$ / Interior) & L5/S1 Compression [BW] \\
\hline
Address (0.211\,s) & 1.344 & 84 / 348 / 14 & 4.87 \\
Top (1.103\,s)     & 0.146 & 152 / 278 / 16 & 9.15 \\
Impact (1.319\,s)  & 0.425 & 145 / 296 / 5  & 5.70 \\
\hline
\end{tabular}
\end{center}

Reserve share ($\sum|r|/\sum|\tau|$) and activations are identical for $p=2$ and
$p=3$ (activation difference maximum $4.7\times 10^{-8}$; reserve difference
$8.4\times 10^{-9}$\,N$\cdot$m) because stage one fixes the reserves and leaves
stage two almost no freedom: only 14, 16, and 5 activations out of 446 are
strictly interior. The cost exponent $p$ is therefore not identifiable on this input.

An address anomaly was identified (filed as \#13427): required moments are small
($|M|\le 35$\,N$\cdot$m per level/axis), yet 49 of 51 DOFs exceed the 1\,N$\cdot$m
reserve limit, axial-rotation reserves ($-18$ to $-54$\,N$\cdot$m) exceed the
axial-rotation demand ($-9$ to $-19$\,N$\cdot$m), and some flexion reserves oppose
the demand (L5--S1 demand $-29.5$\,N$\cdot$m, reserve $+35.2$\,N$\cdot$m).
L5/S1 compression at address of 4.87\,BW reflects co-contraction rather than
demand. Hypotheses include passive fibre force after the Millard swap (affine
model intercept from $a=0.5$ and $a=1$), affine extrapolation error, or thoracic
AR moment-arm screening.

In joint reactions (compression defined as $+F_y$ on the superior by the inferior
in the inferior frame, under \texttt{bruno2015\_regional}, $p=3$), L5/S1
compression is 4.87\,BW at address, 9.15\,BW at top, and 5.70\,BW at impact;
impact peak compression is 8.29\,BW at L2--L3. Maximum resultant shear is
2.11\,BW at address, 2.75\,BW at top, and 8.73\,BW at impact. Context-band ratios
against literature (model peak divided by literature band, descriptive only,
never acceptance) are: Lim 2012 L4/L5 compression peak 1.41; Hosea 1990 L3/L4
professional 0.80, amateur 1.00; Lim 2012 L4/L5 AP shear 4.8, ML shear 6.93.

\paragraph{Limitations and Deferred Work.}
The rigid-tendon affine force model is a declared approximation; activations
below the minimum activation are a model of $a=0$, not a simulated force. The
native run on the Bruno full-body model was executed on capture-A (reported
above), is unqualified, and the address anomaly is open in \#13427. Muscle
wrapping and via-point discontinuities are inherited from the model.

\paragraph{Reproduction.}
\begin{verbatim}
python -m pytest tests/unit/engines/opensim/test_static_opt_core.py \
  tests/opensim/spine/test_static_optimization.py \
  tests/opensim/spine/test_static_optimization_early_samples.py -q -n 0
UD_SPINE_BRUNO_SOURCE=<local MaleFullBody.osim> \
  python3 -m src.engines.physics_engines.opensim.python.spine.study --request base_p2.json
\end{verbatim}
with \texttt{base\_p2.json} holding the request fields listed above
(\texttt{cost\_exponent} 2, \texttt{coupling\_law\_id}
\texttt{bruno2015\_regional}).
